\begin{proof}
\begin{enumerate}[label=\arabic*.]
\item The integral calculation in \cref{maps:src:F.2.6} gives
$\pi_1(M)=0$, $H_2(M;\mathbb Z)=H_3(M;\mathbb Z)=0$, and
torsion-free homology of ranks $1,4,1$ in degrees $0,4,8$.
Hurewicz applied successively gives $\pi_2(M)=\pi_3(M)=0$, and
$e(M)=6$. The complex orientation is the one in
\cref{maps:src:F.3.7}. The characteristic calculations of \cref{night:P1.13} give
\[
Q=I_{3,1},\qquad c_2^2=2,\qquad
p:=p_1/2=-c_2,\qquad p^2=2.
\]
Here $c_1=c_3=0$ because the corresponding integral groups vanish.
The spin Wu identity makes $p$ characteristic for $Q$. Evaluating
its characteristic congruence on each diagonal basis vector shows
that all four coordinates of $p$ are odd.

\item Put $v=(1,1,1,1)^t$. The actual integral reflection is
\[
R=I-v(v^tQ)=
\begin{pmatrix}
0&-1&-1&1\\
-1&0&-1&1\\
-1&-1&0&1\\
-1&-1&-1&2
\end{pmatrix}.
\]
Since $v^tQv=2$, direct multiplication gives $R^2=I$ and
$R^tQR=Q$. Independent coordinate sign changes are also integral
isometries. Use them to write the current vector as $(a,b,c,d)$,
with positive odd entries and
$a^2+b^2+c^2=d^2+2$. For $d>1$, set $s=a+b+c$. Then
\[
s^2-d^2=2(ab+ac+bc)+2>0,\qquad
s^2\leq3(d^2+2)<9d^2.
\]
Thus $d<s<3d$. Reflection subtracts the even integer $s-d$ from
each coordinate; the new last coordinate is $2d-s$, a nonzero odd
integer of absolute value less than $d$. Make the signs positive
again and repeat. The last coordinate strictly decreases until
$d=1$, when the norm equation forces $a=b=c=1$. The resulting
product $U$ of the displayed reflections and sign matrices satisfies
$Up=v$. Its inverse is the reversed product of these involutions.

\item We compute the complete normal clutching datum in fixed
coordinates. Orient $\mathbb H$ by $(1,i,j,k)$ and use
\[
\operatorname{Spin}(4)=S^3_L\times S^3_R,
\qquad (u,w)h=uhw^{-1}.
\]
The covering induces an isomorphism on $\pi_3$. Let $(a,b)$ be
the degrees of the two clutching maps. Evaluation on $1$ gives
$uw^{-1}$, of degree $a-b$, so $e=a-b$, with the sphere orientation
fixed by the left degree-one Euler generator. The left action is
complex linear for right multiplication by $i$, whose real complex
basis $(1,i,j,-k)$ has the opposite quaternion orientation. Its
degree-one generator therefore has $c_2=-1$, and hence $p_1=2$.
The right inverse action is complex linear for left multiplication
by $i$, with real basis $(1,i,j,k)$. Its Euler number and $c_2$ are
both $-1$, so its $p_1$ is again $2$. The determinant-one complex
actions have $c_1=0$. Consequently the exact homomorphism and its
inverse on its image are
\[
(a,b)\longmapsto(e,p_1/2)=(a-b,a+b),
\qquad
(a,b)=\left(\frac{e+p_1/2}{2},\frac{p_1/2-e}{2}\right).
\]
The image consists of pairs of the same parity, and the map is
injective. The tangent clutching class of $S^4$ has $(e,p_1)=(2,0)$,
hence coordinates $(1,-1)$. The homotopy sequence of
$SO(4)\to SO(5)\to S^4$ identifies its span with the kernel of
stabilization and gives surjectivity on $\pi_3$. Subsequent
stabilizations are isomorphisms there, since
$\pi_3(S^n)=\pi_4(S^n)=0$ for $n\geq5$.
The stable coordinate is therefore $a+b=p_1/2$, with the left
generator sent to $1$.

\item For $x\in H_4(M;\mathbb Z)$, use the embedded sphere
representative in the setup of
\cite[\S3.3, pp.~11--12]{IVVWallForm}. Orient its normal bundle
$\nu_x$ by the ambient orientation. Self-intersection and the
tangent splitting give
\[
e(\nu_x)=x\cdot x,\qquad
\langle p_1(\nu_x),[S^4]\rangle=\langle p_1(TM),x\rangle,
\]
because $TS^4$ is stably trivial. Thus its two clutching coordinates
are exactly
\[
\left(\frac{x\cdot x+\langle p,x\rangle}{2},
      \frac{\langle p,x\rangle-x\cdot x}{2}\right),
\qquad S\alpha_M(x)=\langle p,x\rangle.
\]
Their integrality is the characteristic congruence in step~1.
This is the stable normal homomorphism of
\cite[Definition~3.6, p.~12]{IVVWallForm} with its actual generator.
For an oriented middle-cohomology isometry $\psi$ from $M$ to a
comparison manifold $N$, the corresponding homology map is
\[
\Phi=\operatorname{PD}_N\circ\psi\circ\operatorname{PD}_M^{-1}.
\]
When $\psi(p_M)=p_N$, the cup-product evaluation identity gives
$\langle p_N,\Phi(x)\rangle=\langle p_M,x\rangle$.
The map $\Phi$ preserves the intersection pairing and both displayed
normal coordinates, so it intertwines the complete extended forms.

\item Take $N=3\mathbb H\mathbb P^2\#(-\mathbb H\mathbb P^2)$.
On each summand choose the generator $u$ with $p_1=2u$, as in
\cite[Lemma~2.1, p.~5]{ffMandal}. Reversing orientation changes
the evaluation of $u^2$ and leaves the tangent Pontryagin class
unchanged. The collapse maps of the oriented connected sum give
\[
H^4(N;\mathbb Z)=\bigoplus_{i=1}^4\mathbb Z u_i,
\qquad Q_N=I_{3,1},\qquad p_N=u_1+u_2+u_3+u_4;
\]
the tangent formula is
\cite[Theorem~2.4, p.~7]{ffMandal}.
Van Kampen, the connected-sum homology sequence and Hurewicz show
that $N$ is three-connected. The isometry $U$ of step~2 is an
isometry $\psi:H^4(M;\mathbb Z)\to H^4(N;\mathbb Z)$ carrying
$p_M$ to $p_N$. Step~4 supplies the isomorphism of their smooth
extended forms. Wall's classification in
\cite[Theorem~3.7(1), pp.~13--14]{IVVWallForm} now gives an
orientation-preserving diffeomorphism
$M\cong N\#\Sigma$ for an oriented homotopy eight-sphere $\Sigma$.

\item Kramer--Stolz's
\cite[Theorem~A, p.~178]{IVVKramerStolz} applies to
$\mathbb H\mathbb P^2$: it is smooth, simply connected, has dimension
eight, and has total integral homology $\mathbb Z\oplus\mathbb Z
\oplus\mathbb Z$. It gives
$\mathbb H\mathbb P^2\#\Sigma\cong\mathbb H\mathbb P^2$.
This diffeomorphism preserves orientation: its map on $H^4$ sends
a generator to that generator or its negative, whose square still
evaluates to $1$, forcing its degree to be $+1$.
The oriented associativity and permutation maps for connected sum
place $\Sigma$ beside one positive summand. Choose the punctured
balls compatibly with that diffeomorphism and use the identity on
the remaining punctured summands. These maps give
\[
\begin{aligned}
N\#\Sigma
&\cong(\mathbb H\mathbb P^2\#\Sigma)\#\mathbb H\mathbb P^2
       \#\mathbb H\mathbb P^2\#(-\mathbb H\mathbb P^2)\\
&\cong N.
\end{aligned}
\]
Composing with step~5 supplies the required oriented smooth
diffeomorphism. Transporting the integrable structure of
\cref{maps:src:F.3.7} along it gives the asserted complex structure.
\end{enumerate}
\end{proof}
